\chapter{1996年全国卷（理）}
\section{单选题}

\begin{problem}
已知全集 \(I = \N\)，集合 \(A = \{x \mid x = 2n, n \in \N\}\)，\(B = \{x \mid x = 4n, n \in \N\}\)，则（\quad）

\choices
    {\(I = A\cup B\)}
    {\(I = \overline{A}\cup B\)}
    {\(I = A\cup \overline{B}\)}
    {\(I = \overline{A}\cup \overline{B}\)}
\end{problem}

\begin{answer}
C
\end{answer}

\begin{solution}
因为 \(B\subset A\)，对任意 \(x\in I\)，若 \(x\in B\)，则 \(x\in A\)；若 \(x\notin B\)，则 \(x\in\overline B\). 因此 \(I=A\cup\overline B\)，故选 C.
\end{solution}

\begin{problem}
当 \(a > 1\) 时，在同一坐标系中，函数 \(y = a^{-x}\) 与 \(y = \log_a x\) 的图象（\quad）

\choices
    {\choicebitmap[width=0.15\paperwidth]{img/1996/national_science/q2_optA_clean.png}}
    {\choicebitmap[width=0.15\paperwidth]{img/1996/national_science/q2_optB_clean.png}}
    {\choicebitmap[width=0.15\paperwidth]{img/1996/national_science/q2_optC_clean.png}}
    {\choicebitmap[width=0.15\paperwidth]{img/1996/national_science/q2_optD_clean.png}}
\end{problem}

\begin{answer}
A
\end{answer}

\begin{solution}
因为 \(a>1\)，所以 \(y=a^{-x}\) 在 \(\R\) 上单调递减，并经过点 \((0,1)\)；\(y=\log_a x\) 在 \((0,+\infty)\) 上单调递增，并经过点 \((1,0)\)，且当 \(x\to0^+\) 时函数值趋于 \(-\infty\). 与这两个性质相符的图象是选项 A，故选 A.
\end{solution}

\begin{problem}
若 \(\sin^2 x > \cos^2 x\)，则 \(x\) 的取值范围是（\quad）

\choices
    {\(\left\{x\mid 2k\pi -\frac{3}{4}\pi < x < 2k\pi +\frac{1}{4}\pi ,k\in \Z\right\}\)}
    {\(\left\{x\mid 2k\pi +\frac{1}{4}\pi < x < 2k\pi +\frac{5}{4}\pi ,k\in \Z\right\}\)}
    {\(\left\{x\mid k\pi -\frac{1}{4}\pi < x < k\pi + \frac{1}{4}\pi ,k\in \Z\right\}\)}
    {\(\left\{x\mid k\pi +\frac{1}{4}\pi < x < k\pi + \frac{3}{4}\pi ,k\in \Z\right\}\)}
\end{problem}

\begin{answer}
D
\end{answer}

\begin{solution}
由 \(\sin^2x-\cos^2x=-\cos2x\)，题设等价于 \(\cos2x<0\). 因此存在 \(k\in\Z\)，使
\[
\frac{\pi}{2}+2k\pi<2x<\frac{3\pi}{2}+2k\pi
\]
两边同除以 \(2\)，得
\(k\pi+\frac{\pi}{4}<x<k\pi+\frac{3\pi}{4}\)，\(k\in\Z\).
故选 D.
\end{solution}

\begin{problem}
复数 \(\frac{(2 + 2\i)^4}{\left(1 - \sqrt{3}\i\right)^5}\) 等于（\quad）

\choices
    {\(1 + \sqrt{3}\i\)}
    {\(-1 + \sqrt{3}\i\)}
    {\(1 - \sqrt{3}\i\)}
    {\(-1 - \sqrt{3}\i\)}
\end{problem}

\begin{answer}
B
\end{answer}

\begin{solution}
写成三角形式，有
\[
2+2\i=2\sqrt2\left(\cos\frac{\pi}{4}+\i\sin\frac{\pi}{4}\right)
\]
所以
\[
(2+2\i)^4=(2\sqrt2)^4\left(\cos\pi+\i\sin\pi\right)=-64
\]
又
\[
1-\sqrt3\i=2\left(\cos\left(-\frac{\pi}{3}\right)+\i\sin\left(-\frac{\pi}{3}\right)\right)
\]
从而
\[
(1-\sqrt3\i)^5=32\left(\cos\frac{\pi}{3}+\i\sin\frac{\pi}{3}\right)=16+16\sqrt3\i
\]
于是
\[
\frac{(2+2\i)^4}{(1-\sqrt3\i)^5}=\frac{-64}{16+16\sqrt3\i}=-\frac{4}{1+\sqrt3\i}=-1+\sqrt3\i
\]
故选 B.
\end{solution}

\begin{problem}
如果直线 \(l\)，\(m\) 与平面 \(\alpha\)，\(\beta\)，\(\gamma\) 满足：\(l = \beta \cap \gamma\)，\(l \parallel \alpha\)，\(m \subset \alpha\)，\(m \perp \gamma\)，那么必有（\quad）

\choices
    {\(\alpha \perp \gamma\) 且 \(l \perp m\)}
    {\(\alpha \perp \gamma\) 且 \(m \parallel \beta\)}
    {\(m \parallel \beta\) 且 \(l \perp m\)}
    {\(\alpha \parallel \beta\) 且 \(\alpha \perp \gamma\)}
\end{problem}

\begin{answer}
A
\end{answer}

\begin{solution}
因为 \(m\subset\alpha\) 且 \(m\perp\gamma\)，所以平面 \(\alpha\) 与平面 \(\gamma\) 垂直；又因为 \(l\subset\gamma\)，且 \(m\perp\gamma\)，所以 \(m\perp l\). 因此必有 \(\alpha\perp\gamma\) 且 \(l\perp m\)，故选 A.
\end{solution}

\begin{problem}
当 \(-\frac{\pi}{2} \leqslant x \leqslant \frac{\pi}{2}\)，\(f(x) = \sin x + \sqrt{3} \cos x\) 的（\quad）

\choices
    {最大值是 \(1\)，最小值是 \(-1\)}
    {最大值是 \(1\)，最小值是 \(-\frac{1}{2} \)}
    {最大值是 \(2\)，最小值是 \(-2\)}
    {最大值是 \(2\)，最小值是 \(-1\)}
\end{problem}

\begin{answer}
D
\end{answer}

\begin{solution}
有
\[
f(x)=\sin x+\sqrt3\cos x=2\sin\left(x+\frac{\pi}{3}\right)
\]
当 \(-\frac{\pi}{2}\leqslant x\leqslant\frac{\pi}{2}\) 时，\(-\frac{\pi}{6}\leqslant x+\frac{\pi}{3}\leqslant\frac{5\pi}{6}\). 其中 \(x=\frac{\pi}{6}\) 时正弦值为 \(1\)，故最大值为 \(2\)；区间左端正弦值为 \(-\frac12\)，且区间内没有更小的正弦值，故最小值为 \(-1\). 因此选 D.
\end{solution}

\begin{problem}
椭圆 \(\begin{cases}
x &= 3 + 3\cos\varphi,\\
y &= -1 + 5\sin\varphi
\end{cases}\) 的两个焦点坐标是（\quad）

\choices
    {\((-3,5)\)，\((-3,-3)\)}
    {\((3,3)\)，\((3,-5)\)}
    {\((1,1)\)，\((-7,1)\)}
    {\((7,-1)\)，\((-1,-1)\)}
\end{problem}

\begin{answer}
B
\end{answer}

\begin{solution}
由参数方程可知椭圆中心为 \((3,-1)\)，半长轴为 \(5\)，半短轴为 \(3\). 其焦半距
\[
c=\sqrt{5^2-3^2}=4
\]
长轴沿 \(y\) 轴方向，所以两个焦点为 \((3,-1+4)=(3,3)\) 和 \((3,-1-4)=(3,-5)\)，故选 B.
\end{solution}

\begin{problem}
若 \(0 < \alpha < \frac{\pi}{2}\)，则 \(\arcsin \left[\cos \left(\frac{\pi}{2} + \alpha\right)\right] + \arccos \left[\sin (\pi + \alpha)\right]\) 等于（\quad）

\choices
    {\(\frac{\pi}{2} \)}
    {\(-\frac{\pi}{2} \)}
    {\(\frac{\pi}{2}-2\alpha \)}
    {\(-\frac{\pi}{2}-2\alpha \)}
\end{problem}

\begin{answer}
A
\end{answer}

\begin{solution}
因为 \(0<\alpha<\frac{\pi}{2}\)，所以
\(\cos\left(\frac{\pi}{2}+\alpha\right)=-\sin\alpha\)，\(\sin(\pi+\alpha)=-\sin\alpha\).
又 \(-\alpha\in[-\frac{\pi}{2},\frac{\pi}{2}]\)，而 \(\frac{\pi}{2}+\alpha\in[0,\pi]\)，故
\(\arcsin(-\sin\alpha)=-\alpha\)，\(qquad \arccos(-\sin\alpha)=\frac{\pi}{2}+\alpha\).
两式相加得 \(\frac{\pi}{2}\)，故选 A.
\end{solution}

\begin{problem}
将边长为 \(a\) 的正方形 \(ABCD\) 沿对角线 \(AC\) 折起，使得 \(BD = a\)，则三棱锥 \(D - ABC\) 的体积为（\quad）

\choices
    {\(\frac{a^3}{6}\)}
    {\(\frac{a^3}{12}\)}
    {\(\frac{\sqrt{3} a^3}{12}\)}
    {\(\frac{\sqrt{2}}{12} a^3\)}
\end{problem}

\begin{answer}
D
\end{answer}

\begin{solution}
设 \(O\) 为正方形的中心，也是对角线 \(AC\) 的中点. 折叠后 \(OB=OD=\frac{a}{\sqrt2}\)，且 \(AC\) 仍是折叠的公共棱. 在三角形 \(BOD\) 中，由 \(BD=a\) 得
\[
a^2=BD^2=OB^2+OD^2-2OB\cdot OD\cos\angle BOD
=a^2\left(1-\cos\angle BOD\right)
\]
所以 \(\angle BOD=\frac{\pi}{2}\). 又 \(OD\perp AC\)，故 \(OD\) 垂直于平面 \(ABC\) 内相交的两条直线 \(AC\)，\(OB\)，从而 \(OD\) 是三棱锥 \(D-ABC\) 的高. 因此
\[
V_{D-ABC}=\frac13\cdot\frac{a^2}{2}\cdot\frac{a}{\sqrt2}
=\frac{\sqrt2}{12}a^3
\]
故选 D.
\end{solution}

\begin{problem}
等比数列 \(\{a_{n}\}\) 的首项 \(a_1 = -1\)，前 \(n\) 项和为 \(S_{n}\)，若 \(\frac{S_{10}}{S_5} = \frac{31}{32}\)，则 \(\displaystyle\lim_{n\to \infty} S_n\) 等于（\quad）

\choices
    {\(\frac{2}{3}\)}
    {\(-\frac{2}{3}\)}
    {\(2\)}
    {\(-2\)}
\end{problem}

\begin{answer}
B
\end{answer}

\begin{solution}
设等比数列的公比为 \(q\). 若 \(q=1\)，则 \(S_{10}/S_5=2\)，与题设不符，所以 \(q\neq1\). 由等比数列求和公式，
\[
\frac{S_{10}}{S_5}
=\frac{a_1(1-q^{10})/(1-q)}{a_1(1-q^5)/(1-q)}
=1+q^5=\frac{31}{32}
\]
故 \(q^5=-\frac1{32}\)，从而 \(q=-\frac12\). 因为 \(\lvert q\rvert<1\)，
\[
\lim_{n\to\infty}S_n=\frac{a_1}{1-q}
=\frac{-1}{1+\frac12}=-\frac23
\]
故选 B.
\end{solution}

\begin{problem}
椭圆的极坐标方程为 \(\rho = \frac{3}{2 - \cos\theta}\)，则它在短轴上的两个顶点的极坐标是（\quad）

\choices
    {\((3,0)\)，\((1,\pi)\)}
    {\(\left(\sqrt{3},\frac{\pi}{2}\right)\)，\(\left(\sqrt{3},\frac{3\pi}{2}\right)\)}
    {\(\left(2,\frac{\pi}{3}\right)\)，\(\left(2,\frac{5\pi}{3}\right)\)}
    {\(\left(\sqrt{7}, \arctan \frac{\sqrt{3}}{2}\right)\)，\(\left(\sqrt{7}, 2\pi - \arctan \frac{\sqrt{3}}{2}\right)\)}
\end{problem}

\begin{answer}
C
\end{answer}

\begin{solution}
由极坐标关系 \(x=\rho\cos\theta\)，\(x^2+y^2=\rho^2\)，原方程化为
\[
2\rho=3+\rho\cos\theta=3+x
\]
两边平方得
\[
4(x^2+y^2)=(x+3)^2
\]
即
\[
\frac{(x-1)^2}{4}+\frac{y^2}{3}=1
\]
因此椭圆中心为 \((1,0)\)，短半轴为 \(\sqrt3\)，短轴端点为 \((1,\sqrt3)\) 和 \((1,-\sqrt3)\). 它们到极点的距离都是 \(2\)，对应极角分别为 \(\frac{\pi}{3}\) 和 \(\frac{5\pi}{3}\). 故选 C.
\end{solution}

\begin{problem}
等差数列 \(\{a_{n}\}\) 的前 \(m\) 项和为 \(30\)，前 \(2m\) 项和为 \(100\)，则它的前 \(3m\) 项和为（\quad）

\choices
    {\(130\)}
    {\(170\)}
    {\(210\)}
    {\(260\)}
\end{problem}

\begin{answer}
C
\end{answer}

\begin{solution}
设等差数列的公差为 \(d\). 将前 \(2m\) 项分成前后两个各含 \(m\) 项的部分，后 \(m\) 项的每一项都比前面对应项多 \(md\)，所以
\[
S_{2m}=2S_m+m^2d
\]
代入 \(S_m=30\)，\(S_{2m}=100\)，得 \(m^2d=40\). 再将前 \(3m\) 项分成三个各含 \(m\) 项的部分，后三段相对于第一段分别多 \(m^2d\) 和 \(2m^2d\)，故
\[
S_{3m}=3S_m+3m^2d=3\times30+3\times40=210
\]
故选 C.
\end{solution}

\begin{problem}
设双曲线 \(\frac{x^2}{a^2} - \frac{y^2}{b^2} = 1\)（\(0 < a < b\)）的半焦距为 \(c\)，直线 \(l\) 过 \((a,0)\)，\((0,b)\) 两点，已知原点到直线 \(l\) 的距离为 \(\frac{\sqrt{3}}{4} c\)，则双曲线的离心率为（\quad）

\choices
    {\(2\)}
    {\(\sqrt{3}\)}
    {\(\sqrt{2}\)}
    {\(\frac{2\sqrt{3}}{3}\)}
\end{problem}

\begin{answer}
A
\end{answer}

\begin{solution}
直线 \(l\) 的截距式为 \(\frac{x}{a}+\frac{y}{b}=1\)，故原点到 \(l\) 的距离为
\[
\frac{ab}{\sqrt{a^2+b^2}}=\frac{ab}{c}
\]
又双曲线的半焦距满足 \(c^2=a^2+b^2\)，由题设得
\(\frac{ab}{c}=\frac{\sqrt3}{4}c\)，\(\frac{ab}{a^2+b^2}=\frac{\sqrt3}{4}\).
令 \(t=\frac{b}{a}>1\)，则
\[
\frac{t}{1+t^2}=\frac{\sqrt3}{4}
\]
即
\[
\sqrt3t^2-4t+\sqrt3=0
\]
解得 \(t=\sqrt3\) 或 \(t=\frac{\sqrt3}{3}\). 因 \(t>1\)，所以 \(t=\sqrt3\). 因而
\[
e=\frac ca=\sqrt{1+\frac{b^2}{a^2}}=\sqrt{1+3}=2
\]
故选 A.
\end{solution}

\begin{problem}
母线长为 \(1\) 的圆锥体积最大时，其侧面展开图圆心角 \(\varphi\) 等于（\quad）

\choices
    {\(\frac{2\sqrt{2}}{3}\pi\)}
    {\(\frac{2\sqrt{3}}{3}\pi\)}
    {\(\sqrt{2}\pi\)}
    {\(\frac{2\sqrt{6}}{3}\pi\)}
\end{problem}

\begin{answer}
D
\end{answer}

\begin{solution}
设圆锥底面半径为 \(r\)，高为 \(h\). 母线长为 \(1\)，所以 \(h=\sqrt{1-r^2}\)，其中 \(0<r<1\). 体积为
\[
V=\pi r^2\sqrt{1-r^2}
\]
令 \(u=r^2\)，则 \(0<u<1\)，且
\(V(u)=\pi u\sqrt{1-u}\)，\(V'(u)=\frac{\pi(2-3u)}{2\sqrt{1-u}}\).
因此 \(V\) 在 \(u=\frac23\) 时取得最大值，即 \(r=\sqrt{\frac23}\). 侧面展开图是半径为 \(1\) 的扇形，其弧长等于底面周长，所以
\[
\varphi=2\pi r=2\pi\sqrt{\frac23}=\frac{2\sqrt6}{3}\pi
\]
故选 D.
\end{solution}

\begin{problem}
设 \(f(x) \) 是 \((- \infty, + \infty)\) 上的奇函数，\(f(x + 2) = -f(x) \)，当 \(0 \leqslant x \leqslant 1 \) 时，\(f(x) = x\)，则 \(f(7.5)\) 等于（\quad）

\choices
    {\(0.5\)}
    {\(-0.5\)}
    {\(1.5\)}
    {\(-1.5\)}
\end{problem}

\begin{answer}
B
\end{answer}

\begin{solution}
由 \(f(x+2)=-f(x)\) 连续使用两次可得 \(f(x+4)=f(x)\)，所以 \(4\) 是一个周期. 因 \(7.5=-0.5+8\)，
\[
f(7.5)=f(-0.5)
\]
又 \(f\) 是奇函数，且 \(0\leqslant0.5\leqslant1\)，故
\[
f(-0.5)=-f(0.5)=-0.5
\]
因此选 B.
\end{solution}

\section{填空题}

\begin{problem}
已知圆 \(x^{2} + y^{2} - 6x - 7 = 0\) 与抛物线 \(y^{2} = 2px\)（\(p > 0\)）的准线相切，则 \(p =\)\fillinblank{}.
\end{problem}

\begin{answer}
\(2\)
\end{answer}

\begin{solution}
圆的方程可写为
\[
(x-3)^2+y^2=16
\]
所以圆心为 \((3,0)\)，半径为 \(4\)，其左、右端点的横坐标分别为 \(-1\) 和 \(7\). 抛物线 \(y^2=2px\) 的准线为 \(x=-\frac p2\). 由于 \(p>0\)，准线在 \(y\) 轴左侧，只能与圆在左端点处相切，故
\[
-\frac p2=-1
\]
从而 \(p=2\).
\end{solution}

\begin{problem}
正六边形的中心和顶点共 \(7\) 个点，以其中 \(3\) 个点为顶点的三角形共有\fillinblank{}个. （用数字作答）
\end{problem}

\begin{answer}
\(32\)
\end{answer}

\begin{solution}
从 \(7\) 个点中任取 \(3\) 个共有
\[
\mathrm C_7^3=35
\]
种选法. 其中不能构成三角形的情形是中心与一对相对顶点共线. 正六边形共有 \(3\) 对相对顶点，因此应减去 \(3\) 种. 于是三角形共有
\[
35-3=32
\]
个.
\end{solution}

\begin{problem}
\(\tan 20^{\circ} + \tan 40^{\circ} + \sqrt{3}\tan 20^{\circ}\tan 40^{\circ}\) 的值是\fillinblank{}.
\end{problem}

\begin{answer}
\(\sqrt{3}\)
\end{answer}

\begin{solution}
由正切的和角公式，
\[
\tan60^{\circ}
=\frac{\tan20^{\circ}+\tan40^{\circ}}
{1-\tan20^{\circ}\tan40^{\circ}}
=\sqrt3
\]
因此
\[
\tan20^{\circ}+\tan40^{\circ}
=\sqrt3\left(1-\tan20^{\circ}\tan40^{\circ}\right)
\]
移项后得
\[
\tan20^{\circ}+\tan40^{\circ}
+\sqrt3\tan20^{\circ}\tan40^{\circ}=\sqrt3
\]
故填 \(\sqrt3\).
\end{solution}

\begin{problem}
如图，正方形 \(ABCD\) 所在平面与正方形 \(ABEF\) 所在平面成 \(60^{\circ}\) 的二面角，则异面直线 \(AD\) 与 \(BF\) 所成角的余弦值是\fillinblank{}.

\bitmapfigure[float=inline,max width=0.38\linewidth]{img/1996/national_science/q19_fig1.png}
\end{problem}

\begin{answer}
\(\frac{\sqrt{2}}{4}\)
\end{answer}

\begin{solution}
设正方形边长为 \(a\)，取 \(AB\) 为公共棱. 建立空间直角坐标系
\(A=(0,0,0)\)，\(B=(a,0,0)\)，\(D=(0,a,0)\).
由于另一正方形所在平面与 \(ABCD\) 所在平面成 \(60^{\circ}\) 的二面角，可取
\[
F=\left(0,\frac a2,\frac{\sqrt3a}{2}\right)
\]
于是
\(\overrightarrow{AD}=(0,a,0)\)，\(\overrightarrow{BF}=\left(-a,\frac a2,\frac{\sqrt3a}{2}\right)\).
故
\(\overrightarrow{AD}\cdot\overrightarrow{BF}=\frac{a^2}{2}\)，\(|AD|=a\)，\(|BF|=\sqrt2a\).
设异面直线 \(AD\) 与 \(BF\) 所成的锐角为 \(\theta\)，则
\[
\cos\theta=\frac{\left|\overrightarrow{AD}\cdot\overrightarrow{BF}\right|}
{|AD|\,|BF|}
=\frac{a^2/2}{a\cdot\sqrt2a}
=\frac{\sqrt2}{4}
\]
故填 \(\frac{\sqrt2}{4}\).
\end{solution}

\section{解答题}

\begin{problem}
解不等式：\(\log_a\left(1 - \frac{1}{x}\right) > 1\).
\end{problem}

\begin{solution}
对数的底数满足 \(a>0\) 且 \(a\neq1\). 首先由真数为正，得
\[
1-\frac1x>0
\]
所以 \(x<0\) 或 \(x>1\).

当 \(a>1\) 时，对数函数单调递增，原不等式等价于
\[
1-\frac1x>a
\]
若 \(x>1\)，左边小于 \(1\)，不可能大于 \(a>1\)；若 \(x<0\)，移项并乘以负数 \(x\) 时不等号反向，得
\[
1-a>\frac1x
\Longleftrightarrow
(1-a)x<1
\Longleftrightarrow
x>\frac1{1-a}
\]
故此时解集为 \(\frac1{1-a}<x<0\).

当 \(0<a<1\) 时，对数函数单调递减，原不等式等价于
\[
0<1-\frac1x<a
\]
由 \(1-\frac1x<a\) 得 \(\frac1x>1-a>0\)，所以 \(x>0\)，再乘以正数 \(x\) 得
\[
1>(1-a)x
\Longleftrightarrow
x<\frac1{1-a}
\]
结合定义域 \(x>1\)，此时解集为 \(1<x<\frac1{1-a}\).

综上，当 \(a>1\) 时，解集为 \(\left(\frac1{1-a},0\right)\)；当 \(0<a<1\) 时，解集为 \(\left(1,\frac1{1-a}\right)\).
\end{solution}

\begin{problem}
已知 \(\triangle ABC\) 的三个内角 \(A\)，\(B\)，\(C\) 满足：\(A + C = 2B\)，\(\frac{1}{\cos A} + \frac{1}{\cos C} = -\frac{\sqrt{2}}{\cos B}\)，求 \(\cos \frac{A - C}{2}\) 的值.
\end{problem}

\begin{solution}
三角形内角和给出 \(A+B+C=\pi\). 联立 \(A+C=2B\)，得 \(3B=\pi\)，所以 \(B=\frac{\pi}{3}\)，且 \(A+C=\frac{2\pi}{3}\).

设 \(t=\frac{A-C}{2}\)，则 \(A=\frac{\pi}{3}+t\)，\(C=\frac{\pi}{3}-t\)，并且 \(-\frac{\pi}{3}<t<\frac{\pi}{3}\). 由题设，
\[
\frac1{\cos A}+\frac1{\cos C}
=-\frac{\sqrt2}{\cos B}
=-2\sqrt2
\]
另一方面，
\[
\cos A+\cos C
=2\cos\frac{A+C}{2}\cos\frac{A-C}{2}
=\cos t
\]
以及
\[
\cos A\cos C
=\frac{\cos(A+C)+\cos(A-C)}2
=\frac{\cos2t-\frac12}{2}
\]
因此
\[
\frac{\cos t}{(\cos2t-\frac12)/2}=-2\sqrt2
\]
令 \(u=\cos t\). 因为 \(-\frac{\pi}{3}<t<\frac{\pi}{3}\)，所以 \(u>0\)，且 \(\cos2t=2u^2-1\). 代入得
\[
\frac{4u}{4u^2-3}=-2\sqrt2
\Longleftrightarrow
4\sqrt2u^2+2u-3\sqrt2=0
\]
因式分解为
\[
(2u-\sqrt2)(2\sqrt2u+3)=0
\]
由于 \(u>0\)，第二因式不为零，故 \(u=\frac{\sqrt2}{2}\). 所以
\[
\cos\frac{A-C}{2}=\frac{\sqrt2}{2}
\]
\end{solution}

\begin{problem}
如图，在正三棱柱 \(ABC-A_{1}B_{1}C_{1}\) 中，\(E \in BB_{1}\)，截面 \(A_{1}EC \perp\) 侧面 \(AC_{1}\).

\begin{enumerate}
\item 求证：\(BE = EB_{1}\)；
\item 若 \(AA_{1} = A_{1}B_{1}\)，求平面 \(A_{1}EC\) 与平面 \(A_{1}B_{1}C_{1}\) 所成二面角（锐角）的度数.
\end{enumerate}

\bitmapfigure[float=inline,max width=0.26\linewidth]{img/1996/national_science/q22_fig1.png}
\FloatBarrier
\end{problem}

\begin{solution}
设正三棱柱的底面边长为 \(s\)，侧棱长为 \(h\)，建立直角坐标系，使
\(A=(0,0,0)\)，\(B=(s,0,0)\)，\(C=\left(\frac s2,\frac{\sqrt3s}{2},0\right)\)
\(A_1=(0,0,h)\)，\(B_1=(s,0,h)\)，\(C_1=\left(\frac s2,\frac{\sqrt3s}{2},h\right)\).
设 \(E=(s,0,t)\)，其中 \(0\leqslant t\leqslant h\). 侧面 \(ACC_1A_1\) 的一个法向量为 \(\bs{n}_1=(\sqrt3,-1,0)\). 平面 \(A_1EC\) 的一个法向量为
\[
\bs{n}_2=\overrightarrow{A_1C}\times\overrightarrow{A_1E}
=\left(\frac{\sqrt3s(t-h)}2,-\frac{s(h+t)}2,-\frac{\sqrt3s^2}2\right)
\]
由于平面 \(A_1EC\perp ACC_1A_1\)，有 \(\bs{n}_1\cdot\bs{n}_2=0\)，即
\[
\frac{3s(t-h)}2+\frac{s(h+t)}2=s(2t-h)=0
\]
所以 \(t=\frac h2\)，从而
\[
BE=t=\frac h2=EB_1
\]

当 \(AA_1=A_1B_1\) 时，\(h=s\). 此时可将 \(\bs{n}_2\) 乘以 \(-\frac4{s^2}\)，得到平面 \(A_1EC\) 的法向量
\[
\bs{n}=(\sqrt3,3,2\sqrt3)
\]
平面 \(A_1B_1C_1\) 的法向量可取 \(\bs{n}_0=(0,0,1)\). 设两平面所成的锐二面角为 \(\theta\)，则
\[
\cos\theta=\frac{|\bs{n}\cdot\bs{n}_0|}{|\bs{n}|\,|\bs{n}_0|}
=\frac{2\sqrt3}{\sqrt{3+9+12}}
=\frac{\sqrt2}{2}
\]
因此 \(\theta=45^{\circ}\).
\end{solution}

\begin{problem}
某地现有耕地 \(10000\text{ 公顷}\)，规划 \(10\) 年后粮食单产比现在增加 \(22\%\)，人均粮食占有量比现在提高 \(10\%\). 如果人口年增长率为 \(1\%\)，那么耕地平均每年至多只能减少多少公顷（精确到 \(1\text{ 公顷}\)）. （粮食单产 \(=\frac{\text{总产量}}{\text{耕地面积}}\)，人均粮食占有量 \(=\frac{\text{总产量}}{\text{总人口数}}\)）
\end{problem}

\begin{solution}
设耕地平均每年减少 \(x\) 公顷，现有粮食单产为 \(M\)，现有人口为 \(P\). 则 \(10\) 年后耕地面积为 \(10000-10x\)，粮食单产为 \(1.22M\)，人口数为 \(1.01^{10}P\).

现有人均粮食占有量为 \(\frac{10000M}{P}\)，规划后要求的人均粮食占有量为 \(1.10\frac{10000M}{P}\). 因而规划后粮食总产量至少满足
\[
1.22M(10000-10x)
\geqslant
1.10\frac{10000M}{P}\cdot1.01^{10}P
\]
约去 \(M>0\)，得
\[
x\leqslant1000\left(1-\frac{1.10\times1.01^{10}}{1.22}\right)
\]
计算 \(1.01^{10}\approx1.104622\)，所以
\[
x\leqslant1000\left(1-\frac{1.10\times1.104622}{1.22}\right)
\approx4.03
\]
因此耕地平均每年至多只能减少 \(4\) 公顷.
\end{solution}

\begin{problem}
已知 \(l_{1}\)，\(l_{2}\) 是过点 \(P(-\sqrt{2},0)\) 的两条互相垂直的直线，且 \(l_{1}\)，\(l_{2}\) 与双曲线 \(y^{2} - x^{2} = 1\) 各有两个交点，分别为 \(A_{1}\)，\(B_{1}\) 和 \(A_{2}\)，\(B_{2}\).

\begin{enumerate}
\item 求 \(l_{1}\) 的斜率 \(k_{1}\) 的取值范围；
\item 若 \(\left|A_{1}B_{1}\right|=\sqrt{5}\left|A_{2}B_{2}\right|\)，求 \(l_{1}\)，\(l_{2}\) 的方程.
\end{enumerate}
\end{problem}

\begin{solution}
\begin{enumerate}
\item 设 \(l_i\) 的斜率为 \(k_i\)，则
\(l_i:y=k_i\left(x+\sqrt2\right)\)，\(k_1k_2=-1\).
令 \(t=x+\sqrt2\)，则直线上的点为 \((t-\sqrt2,k_it)\). 代入双曲线 \(y^2-x^2=1\)，得
\[
(k_i^2-1)t^2+2\sqrt2t-3=0
\]
要有两个不同交点，首先 \(k_i^2\neq1\)，且判别式
\[
\Delta_i=8+12(k_i^2-1)=4(3k_i^2-1)
\]
必须大于 \(0\)，所以 \(k_i^2>\frac13\). 对 \(i=1\)，\(2\) 同时成立，再利用 \(k_2^2=1/k_1^2\)，得到
\(\frac13<k_1^2<3\)，\(k_1^2\neq1\).
因此
\[
k_1\in\left(-\sqrt3,-1\right)\cup\left(-1,-\frac{\sqrt3}{3}\right)
\cup\left(\frac{\sqrt3}{3},1\right)\cup\left(1,\sqrt3\right)
\]

\item 设斜率为 \(k\) 的直线与双曲线的两个交点对应的参数根为 \(t_1\)，\(t_2\). 由根的判别式，
\[
(t_1-t_2)^2=\frac{4(3k^2-1)}{(k^2-1)^2}
\]
两交点的坐标差为 \((t_1-t_2)(1,k)\)，所以弦长满足
\[
L(k)^2=\frac{4(1+k^2)(3k^2-1)}{(k^2-1)^2}
\]
令 \(u=k_1^2\)，则 \(k_2^2=1/u\). 因而
\[
\frac{|A_1B_1|^2}{|A_2B_2|^2}
=\frac{(1+u)(3u-1)}{(u-1)^2}
\cdot\frac{(1-u)^2}{(1+u)(3-u)}
=\frac{3u-1}{3-u}
\]
由 \(|A_1B_1|=\sqrt5|A_2B_2|\)，得
\[
\frac{3u-1}{3-u}=5
\]
解得 \(u=2\). 所以 \(k_1=\sqrt2\) 或 \(k_1=-\sqrt2\).

当 \(k_1=\sqrt2\) 时，
\(l_1:y=\sqrt2\left(x+\sqrt2\right)\)，\(l_2:y=-\frac{\sqrt2}{2}\left(x+\sqrt2\right)\).
当 \(k_1=-\sqrt2\) 时，
\(l_1:y=-\sqrt2\left(x+\sqrt2\right)\)，\(l_2:y=\frac{\sqrt2}{2}\left(x+\sqrt2\right)\).
\end{enumerate}
\end{solution}

\begin{problem}
已知 \(a\)，\(b\)，\(c\) 是实数，函数 \(f(x) = ax^2 + bx + c\)，\(g(x) = ax + b\)，当 \(-1 \leqslant x \leqslant 1\) 时，\(\left|f(x)\right| \leqslant 1\).

\begin{enumerate}
\item 证明：\(\left|c\right| \leqslant 1\)；
\item 证明：当 \(-1 \leqslant x \leqslant 1\) 时，\(\left|g(x)\right| \leqslant 2\)；
\item 设 \(a > 0\)，当 \(-1 \leqslant x \leqslant 1\) 时，\(g(x)\) 的最大值为 \(2\)，求 \(f(x)\).
\end{enumerate}
\end{problem}

\begin{solution}
\begin{enumerate}
\item 取 \(x=0\)，由题设 \(\left|f(0)\right|\leqslant1\). 因 \(f(0)=c\)，所以
\[
\left|c\right|\leqslant1
\]

\item 由 \(f(1)=a+b+c\)，\(f(0)=c\)，得
\[
g(1)=a+b=f(1)-c
\]
同理，由 \(f(-1)=a-b+c\)，得
\[
g(-1)=-a+b=-f(-1)+c
\]
利用 \(\left|f(1)\right|\leqslant1\)、\(\left|f(-1)\right|\leqslant1\) 及第（1）问结论 \(\left|c\right|\leqslant1\)，有
\(\left|g(1)\right|\leqslant\left|f(1)\right|+\left|c\right|\leqslant2\)，\(\left|g(-1)\right|\leqslant\left|f(-1)\right|+\left|c\right|\leqslant2\).
线性函数 \(g(x)\) 在 \([-1,1]\) 上的值介于 \(g(-1)\) 与 \(g(1)\) 之间，所以对任意 \(-1\leqslant x\leqslant1\)，都有 \(\left|g(x)\right|\leqslant2\).

\item 因 \(a>0\)，\(g(x)=ax+b\) 在 \([-1,1]\) 上递增，故最大值在 \(x=1\) 处取得. 由最大值为 \(2\)，得
\[
g(1)=a+b=2
\]
又 \(g(1)=f(1)-f(0)\)，所以 \(f(1)-f(0)=2\). 由于 \(-1\leqslant f(0)\)，\(f(1)\leqslant1\)，只能有
\(f(0)=-1\)，\(f(1)=1\).
于是 \(c=-1\)，且 \(a+b=2\). 题设还给出 \(f(x)\geqslant-1=f(0)\)，所以
\(x(ax+b)=f(x)-f(0)\geqslant0\)，\(-1\leqslant x\leqslant1\).
当 \(x>0\) 时，\(ax+b\geqslant0\)，令 \(x\to0^+\) 得 \(b\geqslant0\)；当 \(x<0\) 时，\(ax+b\leqslant0\)，令 \(x\to0^-\) 得 \(b\leqslant0\). 因而 \(b=0\)，再由 \(a+b=2\) 得 \(a=2\). 所以
\[
f(x)=2x^2-1
\]
\end{enumerate}
\end{solution}
